% 第6章第9节：本章小结。

\section{本章小结}
\label{sec:ls-analysis-summary}
\label{sec:linear-algebra-summary}
\label{sec:matrix-analysis-summary}

本章围绕一个问题展开：同一个线性作用在什么样的基下最简单，从这个最简形式又能读出它的哪些性质？

前五节逐步调整对基的要求，回答前一半。由特征向量组成的基使矩阵成为对角形；特征值是特征多项式$p_{\bm A}(z)=\det(z\bm I-\bm A)$的根，不随基改变，Cayley--Hamilton定理说明$p_{\bm A}(\bm A)=\bm0$。$n$阶矩阵可对角化，当且仅当它有$n$个线性无关的实特征向量，这时$\bm A^k=\bm V\bm\Lambda^k\bm V^{-1}$。要求基标准正交时，能对角化的恰好是对称矩阵，特征值因而是Rayleigh商的极值；一般方阵只能化为实Schur形，正规矩阵化为分块对角形。输入、输出各取一组标准正交基时，任何矩阵都有奇异值分解，它给出谱范数、Frobenius范数和最佳低秩逼近。

后三节用这些最简形式读出三类性质，它们相互关联而不能互相代替。一次作用的误差放大由条件数$\kappa_2=\sigma_{\max}/\sigma_{\min}$衡量；把矩阵本身当作输入时，对称矩阵的特征值与任意矩阵的奇异值移动不超过$\lVert\bm E\rVert_2$，特征子空间的转角由扰动与谱隙之比控制，非正规矩阵的特征值则可能很敏感；反复作用的长期速度由谱半径决定，$\bm A^k\to\bm0$当且仅当$\rho(\bm A)<1$，但非正规矩阵会先放大再衰减。第~\ref{chap:linear-equations-and-optimization}章求解线性方程、最小二乘与矩阵方程时，将用到奇异值分解、条件数、谱、Cayley--Hamilton定理与实Schur形。
